\section{Introduction}
\label{sec:intro}
%------------------------------------------------------------------------------
The proton and neutron, collectively called the nucleon, are the basic
building blocks of visible matter in the universe today. Ever since they were
discovered in laboratories nearly a century ago~\cite{Rutherford:1919fnt,Chadwick:1932ma}, their fundamental
properties have been vigorously explored:
from the determination of the spin through the specific heat of liquid hydrogen~\cite{Dennison},
to the measurement of the magnetic moments~\cite{Stern}, and the extraction of their electromagnetic sizes through elastic electron scattering
~\cite{Hofstadter:1956qs}. The most revealing discovery, however, came from the
electron deep-inelastic scattering (DIS) on the proton and nuclei at Stanford Linear Accelerator
Center (SLAC) in the late 1960s, in which the constituents of the
proton and neutron, quarks (and later gluons), were discovered~\cite{Bloom:1969kc}.
Soon after, quantum chromodynamics (QCD), a
quantum field theory (QFT) based on ``color'' SU(3) gauge symmetry,
was established as the fundamental theory of strong interactions~\cite{Fritzsch:1973pi,Gross:1973id,Politzer:1973fx},
and of the internal structure of the nucleon as well~\cite{Thomas:2001kw}.

During the last fifty years, significant progress has been made
in understanding the nucleon's internal structure in both experiment and theory.
Multiple experimental facilities have been built to study
high-energy collisions involving protons and nuclei, from which
a large amount of experimental data has been accumulated.
Based on the QCD factorization theorems~\cite{Collins:2011zzd}, derived from perturbative QCD
analyses beyond Feynman's parton model~\cite{Feynman:1973xc},
the parton distribution functions (PDFs), which characterize the
longitudinal momentum distributions of quarks and gluons in hadrons moving at
infinite momentum, have been obtained from global fits
to these data~\cite{Gao:2017yyd, Hou:2019efy, Harland-Lang:2014zoa, Ball:2017nwa}. A recent
result of the phenomenological proton PDFs is shown in Fig.~\ref{fig:PDFS} where $x$
is the momentum fraction of the proton carried by partons.
The PDFs provide a comprehensive
description of the quark and gluon content of the nucleon.
On the theoretical frontier, the Euclidean path-integral formalism of QCD, combined with the lattice
regularization and Monte Carlo simulations~\cite{Wilson:1974sk}, has
offered a systematic way of performing {\it ab initio} calculations of non-perturbative strong interactions.
The rapid rise in computational power and development of intelligent numerical
algorithms have made such a lattice QCD approach extremely successful in computing hadron
spectroscopy, the strong coupling, hadronic form factors,
etc., and even scattering phase shifts~\cite{Tanabashi:2018oca,Briceno:2017max,Aoki:2019cca}.

\begin{figure}[htb]	
	\includegraphics[width=0.9\linewidth]{images/intro-PDFCT18_fig01.pdf}
	\caption{Phenomenological parton distributions obtained by the CTEQ-TEA collaboration (CT18) from
		fits to global high-energy scattering data~\cite{Hou:2019efy}, where $0\le x\le 1$ is the fraction of the proton's
		infinite momentum carried in a parton.}
	\label{fig:PDFS}
\end{figure}

Despite these impressive achievements, we have not been able to systematically
explain the partonic structure of the proton
from first principles, or more explicitly,
we have not made fundamental progress in computing
the quark and gluon distributions starting from the QCD Lagrangian (see \sec{intro_other} for a brief
summary). There is actually a good reason behind it: The
standard formulation of parton physics in the textbooks~\cite{Sterman:1994ce,Collins:2011zzd}
is accomplished through the {\it dynamical correlators} of
quark and gluon fields on the light-front (LF) defined by $t-z={\rm const.}$,
which has the important feature of being independent of the proton's momentum.
On the other hand, lattice QCD is formulated in the Euclidean space with imaginary time,
and cannot be used to directly calculate the dynamical correlations that depend on real time. The standard lattice
approach to parton physics has been to calculate the lower moments of parton distributions,
which are matrix elements of local operators~\cite{Lin:2017snn}.
However, the limitations to the first few moments prohibit practitioners
from reproducing reliably the $x$-dependent structure as shown in Fig.~\ref{fig:PDFS},
other than fitting model functional forms.
Over the years, Hamiltonian diagonalization in LF quantization (LFQ)~\cite{Brodsky:1997de} and
Schwinger-Dyson equations~\cite{Maris:2003vk} have been proposed to solve the nucleon
structure as Minkowskian approaches. Although significant advances have been
made phenomenologically, a systematic approximation to calculate the nucleon PDFs is still missing.

A few years ago, some of the present authors proposed
a general approach to calculate $x$-dependent parton distributions based on
Feynman's original idea about partons: They
are the infinite-momentum limit of {\it static properties of the proton}
at large momentum, and therefore are intrinsically
Euclidean quantities accessible through lattice
QCD~\cite{Ji:2013dva,Ji:2013fga,Ji:2014gla}.
According to this, parton physics in an intermediate
range of $x_{\rm min}\sim 0.1 <x<x_{\rm max}\sim 0.9$ can be calculated
from the physical properties of the proton at a moderately-large
momentum, e.g., with a Lorentz boost factor $\gamma=2-5$. The theory
has been named as {\it large-momentum effective theory} (LaMET)
because a rigorous connection between the infinite-momentum frame (IMF) partons
and quarks and gluons at a finite momentum requires
proper account of the ultraviolet (UV) modes with large momentum
in effective field theory (EFT) and systematic power counting.

The basic principle for LaMET comes from an implicit
observation in the naive parton model:
The structure of the proton is approximately
independent of its momentum so long as it is much larger than
a typical strong-interaction scale $\Lambda_{\rm QCD}$, or its mass.
For example, the quark momentum distribution at moderate $x$ in the proton
at $P=|\vec{P}|=5$ GeV is not very different from that
at $P=50$ GeV or $P=5$ TeV. One might call this phenomenon
{\it large-momentum symmetry}, the nature of which is similar to that of the electronic
structure of the hydrogen atom is not sensitive to the proton mass, so long
as it is much larger than that of the electron.
The asymptotic behavior of the proton structure might be
controlled by an expansion in $\Lambda_{\rm QCD}/P$, but a justification would
require a better understanding of the underlying
dynamics. Assuming this, Feynman replaced the protons probed
at large but finite momenta in high-energy scattering with
the one at infinite momentum $P=\infty$,
corresponding to the leading term in the $\Lambda_{\rm QCD}/P$ expansion, and therefore
the idealized concepts of {\rm the proton in the IMF}
and its constituents---{\it partons}---were born.

In QFTs, however, the existence of the $P=\infty$ limit
depends on their UV behavior. In general, the infinite-momentum limit
does not commute with the UV cut-off limit $\Lambda_{\rm UV}\to \infty$.
While the physical limit is $(\Lambda_{\rm UV}\gg P) \to\infty$,
the parton model and subsequent QCD factorization
theorems use $(P\gg\Lambda_{\rm UV})\to\infty$, keeping
all PDFs with the finite support $|x|\le 1$ where negative $x$ is for antiquarks. Thus partons are an idealized concept
which does not exist in the real world.
Fortunately, because of asymptotic freedom, the above differences can be calculated in perturbative QCD.
Therefore, LaMET is an effective theory of partons, which uses the ordinary
field theoretical calculations $(\Lambda_{\rm UV}\gg P)\to\infty$ and
systematically takes into account non-commuting $P\to\infty$ limits through
EFT matching and running and finite $P$ effects by power corrections. Thus, the PDFs defined in the IMF
or on the LF can be accessed at moderate $x$ from the structure
calculations at $P\sim $ a few GeVs.

The first application of LaMET was to the total gluon helicity $\Delta G$ in the
polarized proton, a quantity of significant experimental interest
at the polarized RHIC~\cite{Bunce:2000uv}, but not within
theoretical reach for many years. In~\cite{Ji:2013fga}, we have shown that from a large-momentum
matrix element of the gluon spin operator in a physical gauge,
$\Delta G$ can be obtained through an EFT matching.
Following this success, LaMET
was applied to the collinear quark PDFs~\cite{Ji:2013dva}. This latter application has
generated considerable theoretical as well as numerical activities,
particularly for the flavor non-singlet $u-d$ distributions in the proton
and other hadrons. A general LaMET framework was subsequently introduced in~\cite{Ji:2014gla}.
More recently, the approach has been extended to the gluons as well~\cite{Zhang:2018diq,Li:2018tpe}.
Therefore, the PDFs can now be computed directly in lattice QCD at specific Feynman variable $x$,
without using LFQ. Besides, the partonic landscape of the proton is extremely
rich, and LaMET holds the promise of computing parton physics beyond the collinear PDFs.

In recent years, tremendous progress has been made in formulating new parton observables
for the proton. In particular, two parallel concepts have been developed
in characterizing the transverse structure of the proton. The first is the generalized
parton distributions (GPDs)~\cite{Mueller:1998fv,Ji:1996ek,Radyushkin:1998es}. The GPDs combine the features of the proton's
elastic form factors, which provide the transverse-space density of partons~\cite{Miller:2007uy},
and Feynman PDFs, and interpolate them. Given the joint
longitudinal-momentum and transverse-space distributions,
one can construct the orbital angular momentum (OAM) of partons, among others~\cite{Ji:1996ek}.
In general, the GPDs can be used to generate momentum-dissected transverse space images of the
proton~\cite{Burkardt:2000za}. A new class of experimental processes, deeply-virtual
exclusive processes (DVEP), including deeply-virtual Compton scattering (DVCS) in which the final state
is a diffractive real photon plus a recoiling proton, has been found to
measure them~\cite{Ji:1996ek,Ji:1996nm}.
The second concept is the transverse-momentum-dependent (TMD) PDFs (or TMDPDFs),
in which the parton's transverse momentum is explicit~\cite{Collins:1981uk,Collins:2011zzd}.
Much theoretical progress has been made in recent years regarding their proper definitions, factorizations,
and spin correlations~\cite{Collins:2012uy,Echevarria:2012js,Collins:2017oxh}.
TMDPDFs can be measured in experimental processes by
observing the transverse momentum of the final-state particles.

\begin{figure}[htb]	
	\centering
	\includegraphics[width=0.95\linewidth]{images/intro-EIC_fig02.pdf}
	\caption{A realization of Electron-Ion Collider at BNL (figure credit to BNL),
		which can be used to probe the partonic landscape of the proton.}
	\label{fig:EIC.png}
\end{figure}

Over the years, it has gradually become clear that a dedicated experimental
facility to fully explore the partonic landscape of the proton is required. To meet this requirement,
the US nuclear science community has proposed, to build a high-energy, high-luminosity
Electron-Ion Collider (EIC)~\cite{Geesaman:2015fha}, which has been recently approved by the US
Department of Energy.
The new collider accelerates electrons to 10-30 GeV and ions--- including the proton and heavy
nuclei all the way up to Pb or U--- up to 100 GeV per nucleon, realizing the center-of-mass
collision energy $E_{\rm cm}$ from 40 to 170 GeV.
The corresponding electron energy in fixed-target experiments
would be 100 GeV to 10 TeV.
The beams are polarized, with high-luminosity up to $10^{33-34}$
collisions/(${\rm cm}^2\cdot{\rm s}$), which are critical
for studying exclusive processes such as DVCS.
The kinematic range of the collisions
covers the Bjorken $x_B$ (which coincides with the parton momentum fraction $x$ in the naive parton model to be discussed in the next section) down to sub-$10^{-4}$, and $Q^2$ as high as $10^4$
GeV$^2$. Much of the EIC science has been discussed in
a dedicated study~\cite{Accardi:2012qut}.

Of course, the EIC and lattice QCD efforts will not stop at the precision parton physics of
the proton. More importantly, we need to develop ways or languages to
describe the nucleon as a strongly-coupled relativistic quantum system, in much
the same way as we understand, for example, the quantum Hall effects in
condensed matter physics. Without a deep understanding of the mechanisms
of  strongly-coupled QCD physics, we cannot claim a fundamental
understanding of the structure of the proton and neutron, in particular,
the origin of their mass and spin. This is one of the most challenging
goals facing the standard model of particle and nuclear physics today.

This review is to systematically expose the idea, formalism, and results
of the LaMET approach to parton physics. We do not claim to be entirely complete
because the field is rapidly developing. References in the related fields
are not meant to be complete either, and we apologize for any important omissions.
Closely-related reviews on lattice parton physics can be found in~\cite{Cichy:2018mum,Zhao:2018fyu}.
There have been studies on the effectiveness of LaMET in various
models~\cite{Ji:2018waw,Gamberg:2014zwa,Broniowski:2017wbr,Xu:2018eii,
	Son:2019ghf,Ma:2019agv,Nam:2017gzm,Bhattacharya:2018zxi,Jia:2015pxx,Hobbs:2017xtq,Broniowski:2017gfp,
	Radyushkin:2017gjd,Bhattacharya:2019cme,Kock:2020frx,DelDebbio:2020cbz},
some of which we will mention in the following for illustrative purposes.
There have been also papers questioning the validity of LaMET method~\cite{Carlson:2017gpk,Rossi:2017muf,Rossi:2018zkn} and some got clarified later in the literature~\cite{Briceno:2017cpo,Ji:2017rah,Radyushkin:2018nbf},
We will not discuss them here and interested readers may refer to the above references.
We have used {\it proton} in most places in the text
to emphasize its importance in nuclear and particle physics. However,
the discussions apply equally to the neutron and other hadrons as well.

The plan for the presentation of this review is as follows. In the remainder of the Introduction, we
explain the nature of parton physics as an effective description
of the internal structure of the proton at large momentum, as well as other existing methods
in the literature for solving the parton structure.
In \sec{lamet}, we introduce the LaMET method starting from
momentum renormalization group equation (RGE) of physical observables in a moving hadron,
followed by the matching between momentum distributions and PDFs. We then formulate an EFT
expansion to compute parton physics from theoretical methods suitable for the structure of a
large-momentum proton. In \sec{renorm}, we discuss some important details for
collinear PDFs: renormalization of the nonlocal operators, particularly
power divergences in lattice regularization, and matching to all orders
in perturbation theory. \sec{gpo} is devoted to applications to general collinear parton observables
including GPDs, parton distribution amplitudes
and higher-order parton correlations. We also discuss applications
for the OAM of the partons in a polarized proton.
In \sec{tmd}, we consider the application to
TMDPDFs, a new class of parton observables. We study matching of the quasi-TMDPDFs to the physical ones,
and explore the lattice calculation of the soft function.
Finally, \sec{lattice} summarizes the recent lattice calculations relevant to
the LaMET applications, and the conclusion is given in \sec{conclusion}.
The review is completed with an Appendix
with a list of acronyms and glossaries, as well as notations and conventions.

%------------------------------------------------------------------------------
\subsection{Partons through Infinite-Momentum States}
\label{sec:lc}
%------------------------------------------------------------------------------
Although partons have become a ubiquitous language
for high-energy scattering, their role as
effective degrees of freedom of QCD for describing
the internal structure of the nucleon is less emphasized
in the literature.
In applications within QCD factorization theorems,
they are---following Feynman---objects arising
from the limit of infinite momentum, with the potential
UV divergences regulated and renormalized after the limit.
Thus,  the partons are an idealized concept,
referring to the quark and gluon Fock components
of the nucleon or other hadrons only in the context of
IMF and LF gauge $A^+=(A^0+A^z)/\sqrt 2=0$.
They are in the same category of concepts as
the {\it infinitely-heavy quark} in heavy-quark effective
field theory (HQET)~\cite{Manohar:2000dt}. To motivate LaMET,
it is important to understand this origin and nature of
partons.

Built from the knowledge of electron scattering in
non-relativistic systems (atoms and molecules)~\cite{West:1974ua},
Feynman introduced the {\it naive parton model} to
describe deep-inelastic scattering (DIS) on the proton, and to explain
the observed phenomenon of Bjorken scaling~\cite{Feynman:1969ej,Bjorken:1969ja,Feynman:1973xc}.

\begin{figure}[htb]	
	\includegraphics[width=0.6\linewidth]{images/intro-dis_fig03.pdf}
	\caption{Deep-inelastic scattering in which partons are probed in the proton.}
	\label{fig:DIS}
\end{figure}

Shown in Fig. \ref{fig:DIS} is the DIS process in which a virtual photon with
large momentum $q^\mu$ is absorbed by a proton of momentum $P^\mu $ and mass $M$.
The invariant variables are $Q^2 = -q^\mu q_\mu $ and $P\cdot q = M\nu$, and
Bjorken $x_B=Q^2/(2P\cdot q)$ fixed in the scaling (or Bjorken) limit $Q^2\to \infty$, $P\cdot q\to \infty$. The inclusive DIS cross section can be factored into a product of leptonic and hadronic tensors, where the former is associated with the electromagnetic current of the lepton, while the latter contains all information about the electromagnetic interaction with the target proton.

To learn about the proton structure, it is best to consider
the scattering in the Breit frame where
\begin{align}\label{eq:diskin}
	q^\mu &= (0,0,0,-Q)  \,, \nonumber \\
	P^\mu &=\left(\sqrt{\frac{Q^2}{4x_B^2}+M^2}, 0, 0, \frac{Q}{2x_B}\right) \,,
\end{align}
and the virtual photon has zero energy. The probe is sensitive only to
the spatial structure as in non-relativistic electron scattering.
However, relativity now constrains the proton to move at a large
momentum $P^z= Q/(2x_B)$ with boost factor $\gamma = Q/(2x_BM)$,
which approaches $P^z=\infty$ in the Bjorken limit.

Feynman made intuitive assumptions about
the proton structure and scattering mechanism, without QFT subtleties~\cite{Feynman:1973xc}:
The proton structure at different large $P^z$ should be similar, and can be
approximated by that at $P^z=\infty$, or in the IMF. The interactions between
constituents (partons) are infinitely time-dilated, and
the wave function configurations are frozen. The proton in high-energy scattering
can be seen as being made of non-interacting partons, each with a longitudinal
momentum $xP^z$ with $0< x< 1$.

The internal structure of non-relativistic
systems is independent of their overall momentum. However, relativistic systems
is different as they least experience the Lorentz contraction.
The structures of such systems are inextricably mixed with
the overall motion, and their dependence on the external momentum is a
dynamical problem. On the other hand, if the internal structure
depends on a particular hadron scale $\Lambda_{\rm QCD}$, the protons at all large-momentum
with $P^z\gg \Lambda_{\rm QCD}$ have a similar structure, corresponding to the
$P^z\to \infty$ limit. This means that if $f(k^z,P^z)$
is the constituent momentum-$k^z$ distribution in a proton of momentum $P^z$,
it might be analytical at $P^z=\infty$ and admits Taylor series expansions
in $1/P^z$,
\begin{align}\label{Eq:expansion}
	f(k^z, P^z) = f(x) + f_2(x)(\Lambda_{\rm QCD}/P^z)^2 + ...\,,
\end{align}
where $x=k^z/P^z$. If so, one may find a large-momentum symmetry of
the proton properties up to power corrections ${\cal O}(1/P^z)$ (we omit the upper index $z$ sometimes
for simplicity), and $f(x)$ is the parton distribution.

The above picture can be shown to hold in certain simple
QFT models, where the dynamical frame dependence of wave functions
for composite systems can be studied straightforwardly.
There are many interesting examples of
two-dimensional systems, for which solutions can be found.
One of the much studied cases is the large $N_c$ QCD, also called
't Hooft model~\cite{tHooft:1974pnl}, in which
the bound states have a well-defined large-momentum limit.
The wave functions can be expanded in
$1/P$, with the corrections starting from $(1/P)^2$.
The momenta of the constituents, $k$ and $P-k$, scale in this
limit. When plotted as a function of $x=k/P$, the change in the
wave function with the magnitude of the momentum can be found
in Figs. 8--11 in~\cite{Jia:2017uul}. This is the type of example
in which Feynman's intuition applies.

However, such a intuition fails in many 3+1 dimensional QFTs, such as QCD.
When a bound state travels at increasingly
large momentum, more and more high-momentum modes of a field theory
are needed to build up its internal structure. Lorentz contraction
indicates that the range of constituent momentum important for the
structure also increases. If these high-momentum modes
do not decouple effectively from the low-momentum ones,
large logarithms of the form $\ln P$, will develop in the structural
quantities. Hence a singularity (cut) at $P=\infty$ can
exist in these theories, making $P\to \infty$ limit ill-defined
and the large momentum expansion impossible. This situation
is intimately related to UV properties of the theories, for which the limits of
taking the UV cut-off $\Lambda_{\rm UV}\to \infty$ and $P\to \infty$
do not commute. While the physically-relevant one is
$(\Lambda_{\rm UV}\gg P)\to \infty$, partons
in QCD factorizations
are obtained in the other limit $(\Lambda_{\rm UV}\ll P)\to \infty$ when
the UV divergences are ignored. Thus one can
formally write the parton distribution as
\begin{align}\label{eq:pd}
	f(x) = \int \frac{d\lambda}{2\pi} e^{ix\lambda}\langle P^z=\infty|\psi^\dagger(z) \psi(0)|P^z=\infty\rangle\,,
\end{align}
where $\lambda = \lim_{P^z\to\infty, z\to 0} (zP^z)$,
and $\psi$ is a quantum field.

Historically, the IMF limit
of field theories has been studied first at the level
of diagrammatic rules for perturbation
theory~\cite{Weinberg:1966jm}. It was found that
taking $P\to\infty$ by ignoring the UV divergences considerably simplifies the perturbation
theory rules: Many time-ordered diagrams vanish and only few
have finite contributions. Moreover, scattering in this limit resembles
that in non-relativistic quantum mechanics, and the wave function
description becomes useful.  The Fock states define
the partons which have the proper kinematic support ($0<x<1$).
After the limit is taken, all physical quantities are
now independent of $P$, and large-momentum symmetry is exact
before UV divergences are regulated. Therefore,
{\it it is the ``naive'' limit, $\Lambda_{\rm UV}\ll P\to \infty$,
	that corresponds to Feynman's naive parton model}.

In the standard QCD study of high-energy scattering,
the above concept of partons as effective degrees of freedom
has been used implicitly. The PDFs are defined in terms of the
naive $P=\infty$ limit, and are used to match the experimental cross sections,
resulting in QCD factorization theorems~\cite{Collins:2011zzd}.

\subsection{Partons through Light-Front Correlators}
\label{sec:lc-imf}

In the literature and textbooks, parton distributions are not traditionally represented
in terms of the Euclidean matrix elements as in \eq{pd}.
Rather, they are represented by the so-called LF correlators of
quantum fields (``operator formalism'')~\cite{Brodsky:1997de,Collins:2011zzd}.
A more explicit formulation in terms of collinear quantum fields
and effective lagrangian is made in the soft-collinear effective theory (SCET)~\cite{Bauer:2000yr,Bauer:2001ct,Bauer:2001yt}.

There is a physical way to see that the parton description of high-energy scattering
results in the light-front correlations. Consider DIS in the rest
frame of the proton, where the virtual photon has momentum
\begin{align}
	q^\mu = (\nu, 0, 0,  -\sqrt{\nu^2+2x_B M\nu}) \,.
\end{align}
In the Bjorken limit $\nu\to\infty$, although the invariant mass $Q$ of the photon
goes to infinity, the photon momentum becomes actually light-like
in the sense that it approaches the light front.
Therefore, in inclusive DIS cross section, the separation of the two electromagnetic currents in the hadronic tensor, which is Fourier conjugate to the photon momentum, also approaches the light-cone direction.

Thus, it appears natural that
all the structural physics of the proton in the IMF can also be
expressed in terms of time-dependent LF correlators or
correlations of quantum fields on the LF. Formally, this is simple to see
if one writes
\begin{align}
	|P\to\infty\rangle = U(\Lambda_\infty) |P=0\rangle \,.
\end{align}
The boost operator $U(\Lambda_\infty)$ can be applied to
the static nonlocal operators in the ordinary momentum distributions.
In doing so, all static correlations become light-cone ones.
The boost process is then similar to shifting the Hamiltonian evolution
in quantum mechanics from Schr\"odinger to Heisenberg
picture where time-dependence is now in the operators.

To express light-cone correlations, it is convenient to introduce
two conjugate light-like (or light-cone) vectors,
$p^\mu=(\Lambda,0,0,\Lambda)$ and $n^\mu=(1/2\Lambda,0,0,-1/2\Lambda)$, with the following properties,
$n^2=p^2=0$, and $n\cdot p=1$, where $\Lambda$ is a parameter. Then any four-vector can be expanded
as,
\begin{align}
	k^\mu = k\cdot n p^\mu + k\cdot p n^\mu  + k_\perp^\mu \ .
\end{align}
In particular, the momentum $P^\mu$ of a proton moving in the $z$-direction can be expressed as
\begin{align}
	P^\mu = p^\mu + (M^2/2)\, n^\mu \,,
\end{align}
where $M$ is the proton mass.

Using the above notation, one can express the unpolarized quark distribution
in the proton as~\cite{Collins:2011zzd},
\begin{align}
	q(x) =\frac{1}{2} \int \frac{d\lambda}{2\pi} e^{i\lambda x}
	\langle P|\overline{\psi}(0) \slashed{n} W(0,\lambda n)\psi(\lambda n)|P\rangle_c \ ,
	\label{eq:LCqPDF}
\end{align}
where $\psi$ is the quark field and $W$ is a gauge-link defined as
\begin{align}
	&W(x_2, x_1)= \nonumber\\
	&\hspace{0.5em}{\cal P}\exp[-ig\int_0^1 dt\,(x_2-x_1)_\mu A^\mu(x_1+(x_2-x_1)t)]
\end{align}
to ensure gauge invariance with ${\cal P}$ denoting the path ordering. $c$ indicates
the connected contributions only, and will be suppressed in the rest of this work.
It is a property of gauge theories in which the charge fields
are not gauge-invariant, and the physical distributions must include a beam of collinear gauge
particles. Note that the above expression is true
for any momentum $P$ (a residual momentum symmetry), in particular,
in the rest frame of the nucleon.
The $x$-support of the above distribution is $[-1,1]$. For negative $x$, one
defines the antiquark distribution with $-q(-x)\equiv \bar q(x)$.
The above expression has been more
familiar in the literature than Feynman's original formulation
of PDFs. In the single quark target, one finds
$q(x)=\delta(x-1)$.

To expose the partons in the above equation, one can follow the QCD light-front quantitzation
~\cite{Chang:1968bh,Kogut:1969xa,Drell:1970yt}, suggested by Dirac in 1949~\cite{Dirac:1949cp}.
In LFQ~\cite{Brodsky:1997de}, one defines the LF coordinates,
\begin{align}
	\xi^\pm = (\xi^0\pm \xi^3)/\sqrt{2} \ ,
\end{align}
where $\xi^+$ is the LF ``time'', and $\xi^-$ is the
LF ``spatial coordinate''. And any four-vector $A^\mu$ will be now written
as $(A^+,A^-,\vec{A}_\perp)$.
Dynamical degrees of freedom are defined on the $\xi^+=0$ plane with
arbitrary $\xi^-$ and $\vec{\xi}_\perp$, with conjugate momentum $k^+$
and $\vec{k}_\perp$. Dynamics is generated
by the light-cone Hamiltonian $H_{\rm LC}=P^-$. For a free particle with
three-momentum $(k^+, \vec{k}_\perp)$ and mass $m$, the {on-shell} LF
energy is $k^-=(\vec{k}_\perp^2+m^2)/(2k^+)$.

For QCD, one can define the Dirac matrices
$\gamma^\pm = (\gamma^0\pm\gamma^3)/\sqrt{2}$, and the
projection operators for the quark fields
as $P_\pm = (1/2)\gamma^\mp\gamma^\pm$, so that any $\psi$
can be decomposed into $\psi=\psi_++\psi_-$ with
$\psi_\pm = P_\pm \psi$, where $\psi_+$ is considered as a dynamical degree of freedom.
For the gauge field, $A^+$ is fixed by the LF gauge
$A^+=0$. $A_\perp$ are dynamical degrees of freedom. $\psi_-$ and $A^-$ are dependent variables, which can be expressed in
terms of $\psi_+$ and $A_\perp$ using equations of motion~\cite{Kogut:1969xa}.

The physics of the LF correlations becomes manifest
if one introduces the canonical expansion,
\begin{align}
	& \psi_+(\xi^+=0,\xi^-,\vec{\xi}_\perp)
	= \int \frac{d^2k_\perp}{ (2\pi)^3}
	\frac{dk^+}{ 2k^+}\sum_\sigma \Big[ b_\sigma(k) u({k,\sigma})
	\nonumber \\
	&  \times e^{-i(k^+\xi^--\vec{k}_\perp\cdot\vec{\xi}_\perp)}
	+ d_\sigma^\dagger(k) v({k,\sigma})e^{i(k^+\xi^--\vec{k}_\perp\cdot\vec{\xi}_\perp)}
	\Big],
\end{align}
where $b^\dagger(k)$ and $d^\dagger(k)$ ($b(k)$ and $d(k)$) are quark and antiquark creation (annihilation)
operators, respectively. {$\sigma$ is the light-cone
	helicity of the quarks which can take $+1/2$ or $-1/2$.} Covariant normalization is adopted for the particle
states and the creation and annihilation operators, i.e.,
\begin{align}
	& \big\{b_\sigma(k),b_{\sigma'}^\dagger(k')\big\}=
	\big\{d_\sigma(k),d_{\sigma'}^\dagger(k')\big\} \nonumber \\
	& =(2\pi)^3\delta_{\sigma\sigma'}2k^+\delta(k^+-k'^{+})
	\delta^{(2)}(\vec{k}_\perp-\vec{k}_\perp^\prime)\,.
\end{align}
Substituting the above expansion into \eq{LCqPDF}, one finds the quark distribution as
\begin{align}
	q(x) = \frac{1}{2x} \sum_\sigma \int \frac{d^2{\vec k}_\perp}{(2\pi)^3}
	\langle P|b_\sigma^\dagger(x,\vec{k}_\perp)b_\sigma(x,\vec{k}_\perp)|P\rangle/\langle P|P\rangle
	\label{eq:partonfock}
\end{align}
for $x>0$, and similarly for $x<0$ for which one gets the antiquark distribution.
The factor $1/x$ comes from the normalization of the creation and annihilation operators. The matrix element above should be interpreted as the matrix element in a wave packet state, in the limit of a state of definite momentum~\cite{Collins:2011zzd}.
This way, one recovers the physical meaning of PDFs in the LF correlator (operator) formalism.

\subsection{Other Approaches to Parton Structure}
\label{sec:intro_other}

Calculating the partonic structure of the hadrons from QCD has always
been an important goal in hadronic physics. There have been two main
approaches apart from various phenomenology and models: light-front quantization
and lattice QCD. Here the authors give a very brief review on LFQ and lattice
approaches that are different from the main subject of this review.

Although LFQ explicitly uses the parton degrees of freedom, it has not been very successful
in practical calculations. First of all, LF perturbation theory,
like the standard Hamiltonian perturbation theory, breaks Lorentz
symmetry manifestly and requires a sophisticated renormalization scheme to restore it.
A potential renormalization scheme must deal with the long-range correlations
in the $\xi^-$ direction which require functional dependence on
the renormalization counterterms~\cite{Wilson:1994fk}. Thus LF perturbation theory has not been used
for any calculations beyond one loop, except for the two-loop anomalous magnetic
moment in QED~\cite{Langnau:1992zj}. In fact, the common wisdom
of using dimensional regularization (DR) for the transverse integrals
and cut-off regularization for the longitudinal one has not been proven useful for
multi-loop calculations, although it has been successfully used to
derive the BFKL evolution by Mueller from the quarkonium wave
functions~\cite{Mueller:1993rr}.

The enthusiasm for using LFQ in QCD is not about perturbation theory, but to solve the hadron states.
Discretized LFQ was proposed in~\cite{Pauli:1985ps} to
make practical calculations for the
bound state problems. This non-perturbative method turns out to be successful
for models in 1+1 dimension, such as the Schwinger
model ~\cite{McCartor:1994im,Harada:1995au}, the 1+1 QCD~\cite{Burkardt:1989wy,Srivastava:2000cf},
the 1+1 $\phi^4$ theory~\cite{Harindranath:1987db} and the sine-Gordon model~\cite{Burkardt:1992sz}.
For 3+1 dimensional theories, simple approximations have been considered, like
the Tamm-Dancoff approximation~\cite{Perry:1990mz}. For QCD itself, one again has
to use severe truncations in the number of Fock states. Some recent works
of this type include~\cite{Vary:2009gt,Lan:2019vui,Jia:2019iyk}.
However, to derive a fully-renormalized hamiltonian is difficult and moreover,
there has been no demonstration so far that the Fock-space truncation actually converges~\cite{Wilson:1994fk}.
Therefore a systematic approximation for QCD bound states in LFQ has yet to be found.

Given the rapid development in lattice QCD, it is natural to use it
to compute parton physics. However, simulating real-time evolution directly
is numerically challenging, which runs into the so-called sign problem or more
generally NP-hard problem. Over the years, a number of methods
have been proposed previously to indirectly calculate the PDFs, which
includes well-studied moment methods, hadronic tensor and Compton amplitude method,
coordinate space factorization, etc. These approaches calculate lattice observables that can be related to the PDFs/structure functions through OPE or the dispersion relation, and thus can be used to probe certain information on the partonic structure of hadrons. However, their aims are
mainly to get the lower moments of PDFs and/or
segments of certain coordinate correlations, not directly in parton
degrees of freedom.

The most-adopted approach on the lattice has been to calculate the moments of PDFs
as the matrix elements of local operators~\cite{Kronfeld:1984zv,Martinelli:1987si}.
In the moments approach, one starts with the so-called twist-two operators~\cite{Christ:1972ms},
\begin{align}
	O^{\mu_1...\mu_n}= {\overline \psi}\gamma^{(\mu_1}iD^{\mu_2} ... iD^{\mu_n)} \psi
	- {\rm trace}
	\label{eq:t2opq}
\end{align}
in the quark case, where $(\mu_1...\mu_2)$ indicates that all the indices are symmetrized, the trace terms are those with at least one factor of the metric
tensor $g^{\mu_i \mu_j}$ multiplied by operators of dimension $(n+2)$ with $n-2$ Lorentz indices, etc.
Their matrix elements in the proton state are
\begin{align}
	\langle P|O^{\mu_1...\mu_n}(\mu)|P\rangle =  2a_n(\mu) (P^{\mu_1}\cdots P^{\mu_n} \!-\! {\rm trace}) \,,
	\label{eq:t2mom}
\end{align}
and the PDFs are related to the local matrix elements through
\begin{align}
	{a_n(\mu)} &= { \int^1_{-1} dx x^{n-1} q(x, \mu^2)}\nn\\
	&= \int^1_0 x^{n-1} \big[q(x, \mu^2)+(-1)^n \bar q(x, \mu^2)\big]
	\label{eq:t2sum}
\end{align}
with $n=1,2,...$. The time-dependent correlation for the PDF
in Eq. (\ref{eq:LCqPDF}) is recovered by taking all the
components as $+$ in Eq. (\ref{eq:t2mom}),
\begin{align}
	\langle P|O^{+...+}(\mu)|P\rangle =2 a_n(\mu) P^+\cdots P^+ \,,
\end{align}
and packaging all the moments into a distribution.
Likewise, for the gluon PDF, its moments are again related
to the matrix elements of local operators,
\begin{align}
	O_g^{\mu_1...\mu_n} = -F^{(\mu_1\alpha}iD^{\mu_2}\cdots .iD^{\mu_{n-1}}F^{\mu_n)}_{\ \alpha}\ ,
	\label{eq:t2opg}
\end{align}
with $n=2,4,6,...$.

A large number of lattice QCD calculations of PDF moments have been done
so far with various degrees of control in systematics~\cite{Lin:2017snn},
which include discretization errors,
physical pion mass, finite volume effects, excited state contaminations,
and proper renormalization. Most of the lattice calculations
have been focused on the first and second moments, $\langle x\rangle$
~\cite{Green:2012ud,Alexandrou:2017oeh,Bali:2014gha},
and $\langle x^2\rangle$\cite{Dolgov:2002zm,Deka:2008xr} for the unpolarized distributions,
and the zero-th and first moments, $\langle 1\rangle$~\cite{Gong:2015iir,Alexandrou:2017oeh,Chang:2018uxx,Alexandrou:2019brg},
and $\langle x\rangle$~\cite{Aoki:2010xg,Abdel-Rehim:2015owa} for the polarized distributions.
However, it has been difficult to
calculate higher moments, due to power divergences
and rapid decay in signals. Nonetheless, moment calculations can provide a useful
calibration for any comprehensive lattice approach to PDFs.

To get more information about the PDFs, it was proposed to
calculate the hadronic tensor of DIS in Euclidean space, and analytically continue the result
to Minkowski space~\cite{Liu:1993cv,Liu:1998um,
	Liu:1999ak,Liu:2016djw,Liu:2017lpe,Liu:2020okp}. Since numerical methods for
analytical continuation are known to be difficult for precision control (similar to NP-hard
or sign problem mentioned earlier), the approach
is useful mainly for the nucleon low-lying excitations. It is very challenging
to obtain parton physics this way.

A similar approach called ``operator product expansion (OPE) without OPE'' was suggested in~\cite{Aglietti:1998ur,Martinelli:1998hz}, see also~\cite{Dawson:1997ic,Capitani:1998fe,Capitani:1999fm}.
The point is that the Compton amplitude in the non-dispersive region
can be calculated in the Euclidean space~\cite{Ji:2001wha}. Through dispersion
relation and Taylor-expansion at $\nu = P\cdot q=0$, one can extract
the higher moments of structure functions from the lattice Compton amplitude.
The recent works and references for parton structure from this approach can be found in~\cite{Chambers:2017dov,Hannaford-Gunn:2020pvu,Horsley:2020ltc}.
A similar method has been adopted for Compton amplitude with heavy-light
currents~\cite{Detmold:2005gg}.
This approach has been used to calculate the second moment of pion distribution amplitude~\cite{Detmold:2018kwu,Detmold:2020lev}.

The current-current correlators can also be studied through OPE
in the coordinate space without momentum insertion into the
currents~\cite{Braun:2007wv}. The spatial correlation at small distances
can be used to calculate higher-moments of distribution amplitudes of the mesons.
A number of lattice studies have been performed in~\cite{Braun:2015axa,Bali:2018spj,Bali:2019dqc}.
Similar strategy has been suggested more recently by Qiu et al.~\cite{Ma:2017pxb}
for parton distributions, and has been used in lattice simulations~\cite{Sufian:2019bol,Sufian:2020vzb}.
The pseudo-PDF has been proposed based on the equal-time correlation---or the quasi-PDF in Fourier space---used
in LaMET~\cite{Radyushkin:2017cyf,Radyushkin:2019owq}, and uses a coordinate-space factorization or OPE at small distance
as in~\cite{Braun:2007wv}. Because of its close connection with the quasi-PDF, we
will discuss comparisons of the pseudo-PDF data analysis method with that for the quasi-PDF
in \sec{renorm-coord}.

There have been pioneering studies on moments of the ``quasi'' quark TMDPDFs on lattice ~\cite{Hagler:2009mb,Musch:2010ka,Musch:2011er,Engelhardt:2015xja,Yoon:2017qzo}.
The staple-shaped gauge link operators have been used to connect the quark fields
separated in the spatial direction to simulate the moments of TMDPDF. The ratios of these moments
are presumed independent of the unknown soft function and may be compared with experimental data.
However, a rigorous relation of these constructions to the physical moments of TMDPDFs
had not been investigated before LaMET, particularly the relationship between
large momentum limit and the rapidity cutoff which is an essential ingredient of
TMD physics. Comparison of this approach and LaMET will be made in \sec{quasitmd}.
%------------------------------------------------------------------------------
